
\section{\trSymplecticDiffeomorphism{}的\trRigidity}
\label{sec:symplectic-diffeo-rigidity}

\hzsync{equation}{3}
\hzsync{theorem}{3}
\hzsync{lemma}{2}
\hzsync{proposition}{7}

我们考察 $(\R^{2n},\omega_0)$ 中的一列\trSymplecticDiffeomorphism{} $\psi_j:\R^{2n}\to\R^{2n}$。%
由定义，其一阶\trDerivative{}满足恒等式
\[
{\psi_j'(x)}^{T}J\,\psi_j'(x)=J,\qquad x\in\R^{2n}.
\]
因此，若序列 $\psi_j$ 依 $C^{1}$ \trConverge，则极限 $\psi(x)=\lim\psi_j(x)$%
也是一个\trSymplecticMapping。与之相对照，我们现在假设序列 $\psi_j$%
仅是局部一致地\trConverge{}到一个\trMapping
\[
\psi(x)=\lim_{j\to\infty}\psi_j(x),
\]
它因此是一个连续映射。由于对每个 $x$ 都有 $\det\psi_j'(x)=1$，%
依据积分的变换公式，我们便得到
\begin{equation}\label{eq:limit-measure-preserving}
\int f\bigl(\psi(x)\bigr)\,dx=\int f(x)\,dx
\end{equation}
对一切 $f\in C_c^{\infty}\bigl(\R^{2n}\bigr)$ 成立，从而 $\psi$ 是\trMeasurePreserving{}的。%
现在再假设 $\psi$ 可微，那么显然 $\det\psi'(x)=\pm1$。%
然而，一个令人瞩目的现象是：只要假定 $\psi$ 可微，它甚至是\trSymplectic{}的，
\[
{\psi'(x)}^{T}J\,\psi'(x)=J,
\]
可见，辛的本性在\trTopology{}极限之下得以幸存。

\begin{theorem}\label{thm:symplectic-embedding-limit}
设 $\varphi_j:B(1)\to\bigl(\R^{2n},\omega_0\bigr)$ 是一列\trSymplecticEmbedding，%
它们局部一致地\trConverge{}到一个\trMapping{} $\varphi:B(1)\to\R^{2n}$。%
若 $\varphi$ 在 $x=0$ 处可微，则 $\varphi'(0)=A$ 是一个\trSymplecticMapping，%
即 $A^{*}\omega_0=\omega_0$。
\end{theorem}

我们看到，一般而言，一个\trVolumePreserving{}的\trDiffeomorphism{}%
不可能在 $C^{0}$ \trTopology{}下用\trSymplecticDiffeomorphism{}来\trApproximation。%
在局部使用\textbf{\trDarboux\trChart}（Darboux chart），%
我们便立刻从\cref{thm:symplectic-embedding-limit}得到

\begin{theorem}[\trEliashberg、\trGromov]\label{thm:sympl-diffeo-c0-closed}
\trCompact\trSymplecticManifold{} $(M,\omega)$ 的\trSymplecticDiffeomorphism{}群，%
在 $M$ 的全体\trDiffeomorphism{}所成之群中是 $C^{0}$-\trClosed{}的。
\end{theorem}

七十年代初，M.~\trGromov{}证明了一个非此即彼的论断：\trSymplecticDiffeomorphism{}群%
要么在全体\trDiffeomorphism{}之群中是 $C^{0}$-\trClosed{}的，%
要么其 $C^{0}$ \trClosure{}就是\trVolumePreserving\trDiffeomorphism{}之群。%
\trSymplecticDiffeomorphism{}能够凭借一些在 $C^{0}$ 极限下稳定的整体性质%
与\trVolumePreserving\trDiffeomorphism{}相区分——八十年代初，Y.~\trEliashberg{}%
在其预印本 ``Rigidity of symplectic and contact structure''（1981，文献[78]）中%
宣告了这一点；该文从未以完整形式发表。部分证明见于\trEliashberg{}[71]（1987）。%
在文献[107]中，\trGromov{}运用\trPseudoholomorphic\trCurve{}的技术%
给出了\cref{thm:sympl-diffeo-c0-closed}%
的一个证明。\trEliashberg{}与\trGromov{}二人都是从非\trEmbedding{}结果推出%
$C^{0}$ 稳定性的。在其专著[108]中，\trGromov{}使用了所谓的\trNash--\trMoser{}%
硬\trImplicitFunctionTheorem{}技巧，而\trEliashberg{}[71]（1987）则使用了%
\cref{thm:symplectic-embedding-limit}的一个类似物。循着 I.~\trEkeland{}与%
H.~\trHofer{}在文献[68]中的策略，下面我们将证明：%
\cref{thm:symplectic-embedding-limit}是任意一个\trCapacity{}函数 $c$ 之存在性的%
一个简单推论。

在下文中，方便的做法是把\trCapacity{}\trContinuation{}到 $\R^{2n}$ 的一切子集上。%
为此，我们取一个定义在开子集 $U\subset\R^{2n}$ 上的\trCapacity{}函数 $c$，%
并对任意子集 $A\subset\R^{2n}$ 定义
\[
c(A)=\inf\bigl\{\,c(U)\ \bigm|\ A\subset U\ \text{且}\ U\ \text{开}\,\bigr\}.
\]
那么，\trMonotonicity
\[
A\subset B\ \Longrightarrow\ c(A)\le c(B)
\]
对 $\R^{2n}$ 的一切子集成立。由\trCapacity{}在\trOpenSet{}上的%
\trSymplectic\trInvariance，对定义在 $A$ 的一个开\trNeighborhood{}上的每一个%
\trSymplecticEmbedding{} $\varphi$，我们可以推出\trInvariance
\[
c\bigl(\varphi(A)\bigr)=c(A).
\]

\begin{proof}[\cref*{thm:symplectic-embedding-limit}的证明]
不失一般性，我们在下文中假设 $\varphi(0)=0$。我们首先断言：%
\trLinear\trMapping{} $\varphi'(0)=A$ 是一个\trIsomorphism。事实上，由于 $\varphi$ 在 $0$ 处可微，%
我们有 $\varphi(x)=Ax+O(|x|)$，于是对半径为 $\varepsilon>0$、以 $0$ 为心的%
开\trBall{} $B_{\varepsilon}$，
\[
\frac{m\bigl(\varphi(B_{\varepsilon})\bigr)}{m(B_{\varepsilon})}
\ \longrightarrow\ \bigl|\det A\bigr|\quad\text{当 }\varepsilon\to0.
\]
另一方面，由于诸\trSymplecticDiffeomorphism{} $\varphi_j$ 是\trVolumePreserving{}的，%
且 $\varphi_j\to\varphi$ 一致，我们有 $m\bigl(\varphi(B_{\varepsilon})\bigr)=m(B_{\varepsilon})$，%
从而 $|\det A|=1$，故 $A$ 是一个\trIsomorphism，正如所断言。%
稍后（\cref{lem:capacity-derivative-isomorphism}）我们将看到，在更弱的假设下%
$A$ 也是\trIsomorphism：我们不再要求 $\varphi_j$ 是\trSymplectic{}的，%
而只要求这些映射保持一个给定的\trCapacity{}函数。

接下来我们断言：要证明\cref{thm:symplectic-embedding-limit}，只需证明
\begin{equation}\label{eq:ast-lambda-omega0}
A^{*}\omega_0=\lambda\omega_0\qquad\text{对某个 }\lambda\neq0.
\end{equation}
事实上，借助 $\varphi_j$ 我们还可以考虑\trSymplecticEmbedding{}%
$(\varphi_j,\id):B(1)\times\R^{2n}\to\bigl(\R^{2n}\times\R^{2n},\omega_0\oplus\omega_0\bigr)$，%
从而对其在 $(0,0)$ 处的\trDerivative{}——由 $\bar{A}=(A,1)$ 给出——得出：%
也有 $\bar{A}^{*}(\omega_0\oplus\omega_0)=\mu(\omega_0\oplus\omega_0)$，%
对某个 $\mu\neq0$。另一方面，由 \eqref{eq:ast-lambda-omega0}，%
$\bar{A}^{*}(\omega_0\oplus\omega_0)=(\lambda\omega_0)\oplus\omega_0$，%
于是 $\mu=1=\lambda$，这正是\cref{thm:symplectic-embedding-limit}所需要的\footnote{原书引作「Theorem 4」，实指论证在先的 Theorem 3，此处径改。——译者注}，%
断言得证。为了证明 \eqref{eq:ast-lambda-omega0}，%
我们要用到下述归功于 Y.~\trEliashberg{}[71]的代数引理。

\begin{lemma}\label{lem:eliashberg-algebraic}
设 $A$ 是一个满足 $A^{*}\omega_0\neq\lambda\omega_0$ 的线性\trIsomorphism。%
那么对每个 $a>0$，存在\trSymplecticMatrix{} $U$ 与 $V$，使得 $U^{-1}AV$%
相对于把 $\R^{2n}=\R^{2}\oplus\R^{2n-2}$ 分解为\trSymplectic\trSubspace{}的分解%
具有形状
\[
U^{-1}AV=\left(\begin{array}{cc|c}
a & 0 & 0\\
0 & a & \\
\hline
* &   & *
\end{array}\right).
\]
\end{lemma}

我们暂且搁置引理的证明，先证 $A^{*}\omega_0=\lambda\omega_0$%
对某个 $\lambda\neq0$ 成立。采用反证法，假设 $A^{*}\omega_0\neq\lambda\omega_0$%
并应用引理。在原点的\trNeighborhood{}中定义\trSymplecticMapping{}%
$\psi_j:=U^{-1}\varphi_jV$，我们得出：$\psi_j$ 局部一致地\trConverge{}到%
$\psi:=U^{-1}\varphi V$，且 $\psi'(0)=U^{-1}AV$。在引理中选取合适的常数 $a$，%
便有 $U^{-1}AV\bigl(B(1)\bigr)\subset Z(1/8)$，从而只要 $\varepsilon$ 充分小，%
就有 $\psi\bigl(B(\varepsilon)\bigr)\subset Z(\varepsilon/4)$\footnote{原书于本处及后两处均在 $\psi$ 或 $\psi_j$ 之外又套一层 $U^{-1}(\cdot)V$，如印作「$U^{-1}\psi V\bigl(B(\varepsilon)\bigr)\subset Z(\varepsilon/4)$」（后两处同型）——而 $\psi=U^{-1}\varphi V$、$\psi_j=U^{-1}\varphi_jV$ 均已有定义，套壳使估计尽失依据，三处一并径去。——译者注}。%
由于 $\psi_j\to\psi$ 局部一致，故当 $j$ 充分大且 $\varepsilon$ 充分小时，%
$\psi_j\bigl(B(\varepsilon)\bigr)\subset Z(\varepsilon/2)$。%
由于 $\psi_j$ 是\trSymplectic{}的，%
由\trCapacity{}的\trInvariance{}与\trMonotonicity{}得出%
$c\bigl(\psi_j(B(\varepsilon))\bigr)=c\bigl(B(\varepsilon)\bigr)
\le c\bigl(Z(\varepsilon/2)\bigr)$，这与\trNontriviality{}公理 (A3) 矛盾。%
\eqref{eq:ast-lambda-omega0} 中的论断已获证明，剩下只需证明%
\cref{lem:eliashberg-algebraic}。

\begin{proof}[\cref*{lem:eliashberg-algebraic}的证明]
\let\qed\relax%
设 $B$ 是 $A$ 的\trSymplectic\trAdjoint，满足
\[
\omega_0(Ax,y)=\omega_0(x,By)
\]
对一切 $x$ 与 $y$ 成立，并简记 $\omega=B^{*}\omega_0$。那么 $\omega\neq\lambda\omega_0$；%
利用 $A$ 是\trIsomorphism{}这一事实，这极易验证。我们断言：存在 $x$，%
使得对每个 $\lambda$ 都有 $\omega(x,\cdot)\neq\lambda\,\omega_0(x,\cdot)$。%
采用反证法，假设对每个 $x$ 都存在 $\lambda(x)\in\R$%
满足 $\omega(x,\cdot)=\lambda(x)\omega_0(x,\cdot)$。若 $x\neq0$，%
则由 $\omega_0$ 的\trNondegenerate{}性，存在 $\xi$ 使得 $\omega_0(\xi,x)\neq0$。%
这对 $x$ 的一个\trNeighborhood{} $U(x)$ 中的一切 $y$ 依然成立。于是
\begin{align*}
\lambda(\xi)\,\omega_0(\xi,y)&=\omega(\xi,y)=-\omega(y,\xi)\\
&=-\lambda(y)\,\omega_0(y,\xi)\\
&=\lambda(y)\,\omega_0(\xi,y),
\end{align*}
这蕴含：对 $x$ 的一个\trNeighborhood{}内的 $y$ 有 $\lambda(\xi)=\lambda(y)$。%
由于 $\R^{2n}\setminus\{0\}$ 是\trConnected{}的，且函数 $\lambda(x)$%
在 $\R^{2n}\setminus\{0\}$ 上局部常值，故它是常数。因此，对 $x\neq0$ 有%
$\omega(x,\cdot)=\lambda\,\omega_0(x,\cdot)$，从而对每个 $x$ 都如此。%
这与 $\omega\neq\lambda\omega_0$ 的假设矛盾，断言于是得证。于是，存在 $x$%
使得\trLinear\trMapping{} $\bigl(\omega_0(x,\cdot),\omega(x,\cdot)\bigr):\R^{2n}\to\R^{2}$%
是\trSurjective。对给定的 $a>0$，我们因此能找到 $y$ 满足
\[
\omega_0(x,y)=1\qquad\text{且}\qquad\omega(x,y)=a^{2}.
\]
回顾 $\omega(x,y)=\omega_0(Bx,By)$，我们可以选取两组\trSymplecticBasis{}%
$(e_1,f_1,\dots)$ 与 $(e_1',f_1',\dots)$，使得 $e_1=x$、$f_1=y$，%
而 $e_1'=\frac{1}{a}Bx$、$f_1'=\frac{1}{a}By$。在这两组基下
\[
Be_1=ae_1',\qquad Bf_1=af_1'.
\]
由 $\langle JAx,y\rangle=\langle Jx,By\rangle=\langle B^{T}Jx,y\rangle$%
我们读出 $A=-JB^{T}J$。现在把 $A$ 在新的基下表示成从带基 $(e_1,f_1,\dots)$%
的 $\R^{2n}$ 到带基 $(e_1',f_1',\dots)$ 的 $\R^{2n}$ 的映射，%
我们便得到所需形状的表示 $U^{-1}AV$。%
这两个\trSymplecticMatrix{}由其列向量定义为 $U=[e_1,f_1,\dots]$%
与 $V=[e_1',f_1',\dots]$。
\end{proof}
\end{proof}

我们知道，\trSymplecticDiffeomorphism{}保持\trCapacity。于是，%
\cref{thm:symplectic-embedding-limit}可以由下面这一归功于 I.~\trEkeland{}与%
H.~\trHofer{}[68]的论断推出——就连续映射而言，它甚至更加出人意料。

\begin{theorem}\label{thm:capacity-continuous-limit}
设 $c$ 是一个\trCapacity。设 $\psi_j:B(1)\to\R^{2n}$ 是一列连续映射，满足
\[
c\bigl(\psi_j(E)\bigr)=c(E)
\]
对一切（小的）\trEllipsoid{} $E\subset B(1)$ 成立，且局部一致地\trConverge{}到
\[
\psi(x)=\lim\psi_j(x).
\]
若 $\psi$ 在 $0$ 处可微，则 $\psi'(0)=A$ 要么是\trSymplectic{}的，要么是\trAntisymplectic{}的：
\[
A^{*}\omega_0=\omega_0\qquad\text{或}\qquad A^{*}\omega_0=-\omega_0.
\]
注意，这里并不要求这些映射可逆。
\end{theorem}

为了证明\cref{thm:capacity-continuous-limit}，我们从下述结果着手。

\begin{lemma}\label{lem:capacity-derivative-isomorphism}
设 $c$ 是一个\trCapacity。考虑 $\R^{2n}$ 中一列连续映射 $\varphi_j$，%
它局部一致地\trConverge{}到\trMapping{} $\varphi$。假设对一切开\trEllipsoid{}%
及一切 $j$ 都有 $c\bigl(\varphi_j(E)\bigr)=c(E)$。若 $\varphi'(0)$ 存在，%
则它是一个\trIsomorphism。
\end{lemma}

\begin{proof}
采用反证法，假设 $A$ 不是\trSurjective，%
于是 $A\bigl(\R^{2n}\bigr)$ 包含在一个\textbf{\trHyperplane}（hyperplane）$H$ 中。%
必要时先与一个线性\trSymplecticMapping{}复合，我们便可假设
\begin{equation}\label{eq:hyperplane-reduction}
A\bigl(\R^{2n}\bigr)\subset H=\bigl\{(x,y)\mid x_1=0\bigr\}.
\end{equation}
由
\[
\psi(x,y)=\Bigl(\frac{1}{\alpha}\,x_1,\,x_2,\dots,x_n,\,\alpha\,y_1,\,y_2,\dots,y_n\Bigr)
\]
定义线性\trSymplecticMapping{} $\psi$，我们可以选取 $\alpha>0$ 充分小，使得
\[
\psi A\bigl(B(1)\bigr)\subset B^{2}\Bigl(\frac{1}{16}\Bigr)\times\R^{2n-2}
=Z\Bigl(\frac{1}{16}\Bigr),
\]
其中右端的开二维\trDisc{} $B^{2}$ 位于以 $\{x_1,y_1\}$ 为坐标的%
\trSymplectic{}平面内。按\trDerivative{}的定义，%
我们有 $|\psi\varphi(x)-\psi A(x)|\le a(|x|)|x|$，其中当 $s\to0$ 时 $a(s)\to0$。%
于是，只要 $\varepsilon$ 充分小，就有
\[
\psi\varphi\bigl(B(\varepsilon)\bigr)\subset Z\Bigl(\frac{\varepsilon}{4}\Bigr).
\]
由于 $\psi\varphi_j$ 局部一致地\trConverge{}到 $\psi\varphi$，故当 $j$ 充分大时
\[
\psi\varphi_j\bigl(B(\varepsilon)\bigr)\subset Z\Bigl(\frac{\varepsilon}{2}\Bigr),
\]
由假设 $\psi\varphi_j$ 保持\trCapacity，于是由\trMonotonicity
\[
c\bigl(B(\varepsilon)\bigr)=c\Bigl(\psi\varphi_j\bigl(B(\varepsilon)\bigr)\Bigr)
\le c\Biggl(Z\Bigl(\frac{\varepsilon}{2}\Bigr)\Biggr)
=\frac{1}{4}\,c\bigl(B(\varepsilon)\bigr).
\]
这一矛盾表明 $A$ 是\trSurjective。
\end{proof}

\begin{proof}[\cref*{thm:capacity-continuous-limit}的证明]
不妨设 $\psi(0)=0$。由\cref{lem:capacity-derivative-isomorphism}，$A=\psi'(0)$%
是一个\trIsomorphism，我们首先证明 $A^{*}\omega_0=\lambda\omega_0$%
对某个 $\lambda\neq0$ 成立。采用反证法，假设 $A^{*}\omega_0\neq\lambda\omega_0$，%
并（由\cref{lem:eliashberg-algebraic}）找到满足%
$U^{-1}AV\bigl(B(1)\bigr)\subset Z(1/8)$ 的\trSymplecticMapping{} $U$ 与 $V$。%
现在仿照\cref{thm:symplectic-embedding-limit}的证明继续\footnote{原书引作「Theorem 4」，实指论证在先的 Theorem 3，此处径改。——译者注}，%
定义序列 $\varphi_j:=U^{-1}\psi_jV$。那么%
$\varphi_j\to\varphi:=U^{-1}\psi V$ 局部一致，且 $\varphi'(0)=U^{-1}AV$。%
于是 $\varphi\bigl(B(\varepsilon)\bigr)\subset Z(\varepsilon/4)$，%
从而对充分大的 $j$ 与充分小的 $\varepsilon>0$ 有%
$\varphi_j\bigl(B(\varepsilon)\bigr)\subset Z(\varepsilon/2)$。%
根据对 $\psi_j$ 的假设，$c\bigl(\varphi_j(B(\varepsilon))\bigr)=c\bigl(B(\varepsilon)\bigr)$，%
我们由\trCapacity{}的\trMonotonicity{}推知%
$c\bigl(B(\varepsilon)\bigr)\le c\bigl(Z(\varepsilon/2)\bigr)$，%
这与\trCapacity{}的公理 (A3) 矛盾。

$A^{*}\omega_0=\lambda\omega_0$ 已获证明。由\trConformality，%
线性\trAntisymplectic{}映射保持\trCapacity。把\trMapping{} $\psi_j$ 与 $\psi$ 在 $\lambda>0$ 时%
与\trSymplecticMapping{} $B=\Bigl(\frac{1}{\sqrt{\lambda}}A\Bigr)^{-1}$ 复合、%
在 $\lambda<0$ 时与\trAntisymplectic{}映射 $B=\Bigl(\frac{1}{\sqrt{-\lambda}}A\Bigr)^{-1}$ 复合，%
我们便化归到情形
\[
A=\alpha\,1,\qquad\text{其中 }\alpha>0,
\]
而尚需证明 $\alpha=1$。若 $\alpha<1$，%
则存在一个小\trBall{}与一个满足 $\alpha<r<1$ 的 $r$，%
使得对大的 $j$ 有 $\psi_j\bigl(B(\varepsilon)\bigr)\subset B(r\varepsilon)$。%
由于 $\psi_j$ 保持\trCapacity，我们得到%
$c\bigl(B(\varepsilon)\bigr)=c\bigl(\psi_j(B(\varepsilon))\bigr)
\le c\bigl(B(r\varepsilon)\bigr)=r^{2}c\bigl(B(\varepsilon)\bigr)$，%
这是一个矛盾。在情形 $\alpha>1$ 时我们将证明
\begin{equation}\label{eq:superset-ball}
\psi_j\bigl(B(\varepsilon)\bigr)\supset B(r\varepsilon)
\end{equation}
对某个满足 $\alpha>r>1$ 的 $r$、小的 $\varepsilon$ 与大的 $j$ 成立，%
这将导致矛盾 $r^{2}c\bigl(B(\varepsilon)\bigr)=c\bigl(B(r\varepsilon)\bigr)
\le c\bigl(B(\varepsilon)\bigr)$。从而 $\alpha=1$，正如定理所断言。

为证明 \eqref{eq:superset-ball}，我们必须证明：对每个 $y\in B(r\varepsilon)$，%
存在 $x\in B(\varepsilon)$ 满足 $\psi_j(x)=y$。%
我们采用一个以\trBrouwer\trMapping\trDegree{}为基础的指标论证。固定 $1<r<\alpha$。那么，%
注意到 $A\bigl(B(\varepsilon)\bigr)=B(\alpha\varepsilon)$，对 $y\in B(r\varepsilon)$ 便有
\[
\deg\bigl(B(\varepsilon),A,y\bigr)=1,
\]
而只要我们能证明
\[
\deg\bigl(B(\varepsilon),\psi_j,y\bigr)=\deg\bigl(B(\varepsilon),A,y\bigr),
\]
证明即告完成。依据\trDegree{}的\trHomotopy\trInvariance，只需验证\trHomotopy{}%
$h(t,x)=t\psi_j(x)+(1-t)Ax$ 满足：当 $x\in\partial B(\varepsilon)$%
且 $0\le t\le1$ 时 $h(t,x)\neq y$。

回顾 $\psi(0)=0$，从而 $|\psi(x)-Ax|\le a(|x|)|x|$，%
其中当 $s\to0$ 时 $a(s)\to0$。由于 $\psi_j\to\psi$ 在\trCompact{}集上一致，%
故对每个 $\sigma>0$ 与 $j\ge j_0(\sigma)$，
\begin{equation}\label{eq:uniform-approximation}
|\psi_j(x)-Ax|\le a(|x|)|x|+\sigma.
\end{equation}
采用反证法，假设对 $x\in\partial B(\varepsilon)$ 与 $y\in B(r\varepsilon)$ 有%
$t\psi_j(x)+(1-t)Ax=y$。那么
\[
t\bigl(\psi_j(x)-Ax\bigr)=y-Ax.
\]
右端大于 $|Ax|-|y|=\alpha\varepsilon-|y|\ge(\alpha-r)\varepsilon$；%
而左端依据 \eqref{eq:uniform-approximation} 却小于 $a(|\varepsilon|)\varepsilon+\sigma$。%
于是，取 $\sigma=a(|\varepsilon|)\varepsilon$ 并取 $\varepsilon$ 充分小，%
便得出矛盾，从而证明了：对一切 $0\le t\le1$ 与一切 $x\in\partial B(\varepsilon)$%
都有 $h(t,x)\neq y$。\cref{thm:capacity-continuous-limit}的证明到此完成。
\end{proof}

作为一例有趣的特殊情形，我们从\cref{thm:capacity-continuous-limit}得到下述

\begin{corollary}\label{cor:linear-map-ellipsoid-capacity}
$\R^{2n}$ 中的\trLinear\trMapping{} $A$%
若保持诸\trEllipsoid{}的\trCapacity，即 $c\bigl(A(E)\bigr)=c(E)$，%
则它要么是\trSymplectic{}的，要么是\trAntisymplectic{}的：
\[
A^{*}\omega_0=\omega_0\qquad\text{或}\qquad A^{*}\omega_0=-\omega_0.
\]
\end{corollary}

设 $c$ 是任一个\trCapacity{}函数，并考虑 $\R^{2n}$ 的一个满足
\[
c\bigl(h(E)\bigr)=c(E)\qquad\text{对一切（小的）\trEllipsoid{} $E$}
\]
的\textbf{\trHomeomorphism}（homeomorphism）$h$。那么，若 $h$ 还是可微的，%
我们从\cref{thm:capacity-continuous-limit}得出：$h$%
是一个要么\trSymplectic{}要么\trAntisymplectic{}的\trDiffeomorphism，%
即 $h'(x)^{*}\omega_0=\pm\omega_0$。%
这与\trMeasurePreserving{}的\trHomeomorphism——%
即满足 \eqref{eq:limit-measure-preserving} 的\trHomeomorphism——相仿；%
在后一情形，当 $h$ 可微时我们得出 $\det h'(x)=\pm1$。我们看到，%
每个\trCapacity{}函数 $c$ 都甄别出一个特定的群：%
保持一切\trOpenSet{}之\trCapacity{}的\trHomeomorphism{}所成之群。%
这个\trHomeomorphism{}群的元素还具有如下附加性质：%
只要它们可微，它们就是\trSymplectic{}的或\trAntisymplectic{}的。因此，这个群可以视为%
\trSymplecticDiffeomorphism{}群的一个\trTopology{}版本。%
这个群是否在局部一致极限下封闭，尚属未知。但下述较弱的结果成立。

\begin{theorem}\label{thm:homeomorphism-capacity-limit}
设 $h_j:\R^{2n}\to\R^{2n}$ 是一列\trHomeomorphism，满足
\[
c\bigl(h_j(E)\bigr)=c(E)
\]
对一切（相应地，一切小的）\trEllipsoid{} $E$ 成立。设 $h_j$%
局部一致地\trConverge{}到 $\R^{2n}$ 的一个\trHomeomorphism{} $h$。那么
\[
c\bigl(h(E)\bigr)=c(E)
\]
对一切（相应地，一切小的）\trEllipsoid{} $E$ 成立。
\end{theorem}

\begin{proof}
由于 $h^{-1}\circ h_j\to\id$ 局部一致，%
故对每个\trEllipsoid{} $E$ 与每个 $0<\varepsilon<1$，只要 $j$ 充分大便有
\[
h^{-1}\circ h_j\bigl((1-\varepsilon)E\bigr)\subset E
\subset h^{-1}\circ h_j\bigl((1+\varepsilon)E\bigr),
\]
这一\trTopology{}事实极易借助与上文相同的\trDegree{}论证来验证。于是
\[
h_j\bigl((1-\varepsilon)E\bigr)\subset h(E)\subset h_j\bigl((1+\varepsilon)E\bigr).
\]
由\trMonotonicity{}与\trConformality，%
$(1-\varepsilon)^{2}c(E)=c\bigl((1-\varepsilon)E\bigr)
=c\bigl(h_j((1-\varepsilon)E)\bigr)\le c\bigl(h(E)\bigr)$，%
且 $c\bigl(h(E)\bigr)\le c\bigl(h_j((1+\varepsilon)E)\bigr)
=c\bigl((1+\varepsilon)E\bigr)=(1+\varepsilon)^{2}c(E)$。简言之，
\[
(1-\varepsilon)^{2}c(E)\le c\bigl(h(E)\bigr)\le(1+\varepsilon)^{2}c(E).
\]
这对每个 $\varepsilon>0$ 都成立，故如所愿有 $c\bigl(h(E)\bigr)=c(E)$。
\end{proof}

为了推广这一论断，我们用 $\mathcal{O}$ 表示 $\R^{2n}$ 中%
开且\trBounded{}之集所成的族，并对 $\Omega\in\mathcal{O}$ 配以
\begin{align*}
\check{c}(\Omega)&=\sup\bigl\{\,c(U)\ \bigm|\ U\ \text{开且}\ \overline{U}\subset\Omega\,\bigr\}\\
\hat{c}(\Omega)&=\inf\bigl\{\,c(U)\ \bigm|\ U\ \text{开且}\ U\supset\overline{\Omega}\,\bigr\}.
\end{align*}
\trOpenSet{}的特定族 $\mathcal{O}_{c}$ 由条件
\[
\mathcal{O}_{c}=\bigl\{\Omega\in\mathcal{O}\ \bigm|\ \hat{c}(\Omega)=\check{c}(\Omega)\bigr\}.
\]
来定义。

\begin{proposition}\label{prop:homeomorphism-regular-family}
设 $h_j$ 是 $\R^{2n}$ 的一列\trHomeomorphism，%
局部一致地\trConverge{}到 $\R^{2n}$ 的一个\trHomeomorphism{} $h$。%
若对一切 $U\in\mathcal{O}$ 与一切 $j$ 都有 $c\bigl(h_j(U)\bigr)=c(U)$，%
则 $c\bigl(h(\Omega)\bigr)=c(\Omega)$ 对一切 $\Omega\in\mathcal{O}_{c}$ 成立。
\end{proposition}

\begin{proof}
若 $\Omega\in\mathcal{O}_{c}$ 且 $\varepsilon>0$，%
我们找到 $\hat{U},\check{U}\in\mathcal{O}$，%
满足 $\hat{U}\supset\overline{\Omega}$ 与 $\overline{\check{U}}\subset\Omega$，使得
\[
c(\hat{U})\le c(\Omega)+\varepsilon\qquad\text{且}\qquad c(\check{U})\ge c(\Omega)-\varepsilon.
\]
对大的 $j$，由上述\trDegree{}论证有%
$h_j(\hat{U})\supset h(\Omega)\supset h_j(\check{U})$，%
于是利用\trCapacity{}的\trMonotonicity，
\[
c(\hat{U})=c\bigl(h_j(\hat{U})\bigr)\ge c\bigl(h(\Omega)\bigr)
\ge c\bigl(h_j(\check{U})\bigr)=c(\check{U}),
\]
从而 $c(\Omega)+\varepsilon\ge c\bigl(h(\Omega)\bigr)\ge c(\Omega)-\varepsilon$。%
这对每个 $\varepsilon>0$ 都成立，命题得证。
\end{proof}

\begin{definition}\label{def:inner-outer-regular-capacity}
一个\trCapacity{} $c$ 称为内\trRegular{}的，%
相应地，称为外\trRegular{}的，若
\[
c(U)=\check{c}(U)\qquad\text{相应地，若}\qquad c(U)=\hat{c}(U)
\]
对一切 $U\in\mathcal{O}$ 成立。
\end{definition}

\begin{proposition}\label{prop:regular-capacity-homeomorphism}
设\trCapacity{} $c$ 是内\trRegular{}的或外\trRegular{}的。设 $\varphi_j$ 与 $\varphi$%
是 $\R^{2n}$ 的\trHomeomorphism，且
\[
\varphi_j\to\varphi\qquad\text{且}\qquad\varphi_j^{-1}\to\varphi^{-1}
\]
局部一致。若对一切 $U\in\mathcal{O}$ 与一切 $j$ 都有%
$c\bigl(\varphi_j(U)\bigr)=c(U)$，那么也有
\[
c\bigl(\varphi(U)\bigr)=c(U)\qquad\text{对一切}\quad U\in\mathcal{O}.
\]
\end{proposition}

\begin{proof}
a) 设 $c$ 是内\trRegular{}的，并设 $\Omega$ 开且\trBounded。那么，%
若 $U$ 开且 $\overline{U}\subset\Omega$，则当 $j$ 大时有%
$\varphi_j(U)\subset\varphi(\Omega)$，%
从而 $c\bigl(\varphi_j(U)\bigr)=c(U)\le c\bigl(\varphi(\Omega)\bigr)$，%
于是 $c(U)\le c\bigl(\varphi(\Omega)\bigr)$。因此，取上确界便得%
$c(\Omega)\le c\bigl(\varphi(\Omega)\bigr)$。类似地可证，%
对每个开且\trBounded{}的 $\Omega$ 都有 $c(\Omega)\le c\bigl(\varphi^{-1}(\Omega)\bigr)$。%
于是，由于 $\varphi$ 是\trHomeomorphism，%
$c(\Omega)=c\bigl(\varphi^{-1}\circ\varphi(\Omega)\bigr)\ge c\bigl(\varphi(\Omega)\bigr)
\ge c(\Omega)$，从而 $c\bigl(\varphi(\Omega)\bigr)=c(\Omega)$，正如所求。

b) 若 $c$ 是外\trRegular{}的且 $\Omega$ 开且\trBounded，则对 $\varepsilon>0$ 我们得出：%
当 $j$ 充分大时 $\varphi_j(\Omega)\subset U_{\varepsilon}\bigl(\varphi(\Omega)\bigr)$；%
这里 $U_{\varepsilon}=\{x\mid\operatorname{dist}(x,U)<\varepsilon\}$。于是%
$c(\Omega)=c\bigl(\varphi_j(\Omega)\bigr)\le c\bigl(U_{\varepsilon}(\varphi(\Omega))\bigr)$，%
对右端取下确界便得 $c(\Omega)\le c\bigl(\varphi(\Omega)\bigr)$。%
照 a) 部分那样论证，我们得出 $c\bigl(\varphi(\Omega)\bigr)=c(\Omega)$%
对每个 $\Omega\in\mathcal{O}$ 成立。
\end{proof}

到目前为止，我们已经从\trSymplecticCapacity{}的存在性推出了一些%
关于\trSymplectic{}映射的惊人现象。然而，这个概念本身仍然相当神秘，%
并引出许多问题。例如，小集合之\trCapacity{}的知识是否足以理解较大集合的%
\trCapacity，这一点尚属未知。说得更确切些：比如说，%
一个保持小集合之\trCapacity{}的\trHomeomorphism{}是否保持大集合的\trCapacity？

我们尚没有一个保持某一\trCapacity{}而不保持另一\trCapacity{}的%
\trHomeomorphism{}的例子。我们同样不知道，保持\trOpenSet{}之\trCapacity{}的%
\trHomeomorphism{}是否也保持\trOpenSet{}的\trLebesgue{}\trMeasure。%
但容易看出，保持\trOpenSet{}的一切\trCapacity{}的\trHomeomorphism{}%
必定是\trMeasurePreserving{}的。为了证明这一点，%
我们首先定义一个特殊的\trEmbedding\trCapacity{} $\gamma$。%
对 $r>0$ 引入\textbf{开\trCube}（open cube）
\[
Q(r)=(0,r)^{2n}\subset\R^{2n},
\]
其棱平行于坐标轴，我们定义
\[
\gamma(M,\omega):=\sup\bigl\{\,r^{2}\ \bigm|\ \text{存在\trSymplecticEmbedding{} }\varphi:Q(r)\to M\,\bigr\}.
\]
显然，$\gamma$ 是一个满足公理 (A1)、(A2)%
与弱\trNontriviality{}条件 (A3$'$) 的\trCapacity。它不是规范化的。%
利用\trEkeland{}与\trHofer{}在文献[69]中的 $n$ 阶\trCapacity{}函数 $c_n$，%
可以证明 $\gamma\bigl(B(r)\bigr)=\tfrac{1}{n}\pi r^{2}$；%
而且 $\gamma\bigl(Z(r)\bigr)=\pi r^{2}$。

\begin{proposition}\label{prop:gamma-preserves-measure}
设 $h$ 是 $\R^{2n}$ 的一个\trHomeomorphism，满足
\[
\gamma\bigl(h(\Omega)\bigr)=\gamma(\Omega)
\]
对一切 $\Omega\in\mathcal{O}$ 成立。那么 $h$ 保持\trOpenSet{}的%
\trLebesgue{}\trMeasure{} $\mu$，即
\[
\mu\bigl(h(\Omega)\bigr)=\mu(\Omega)
\]
对一切 $\Omega\in\mathcal{O}$ 成立。
\end{proposition}

\begin{proof}
由 $\gamma$ 的定义我们推得
\[
\mu\bigl(Q(r)\bigr)=r^{2n}=\gamma\bigl(Q(r)\bigr)^{n}.
\]
由于每个\trSymplecticEmbedding{}都是\trVolumePreserving{}的，%
故对开且\trBounded{}集 $\Omega\subset\R^{2n}$ 的\trCapacity{} $\gamma(\Omega)$，%
我们得到
\[
\mu(\Omega)\ \ge\ \sup_{r}\mu\bigl(Q(r)\bigr)\ =\ \gamma(\Omega)^{n},
\]
其中上确界取遍那些存在\trSymplecticEmbedding{} $Q(r)\to\Omega$ 的 $r$。%
若 $Q$ 是任一个棱平行于坐标轴的开\trCube，我们便得出
\[
\mu(Q)=\gamma(Q)^{n}=\gamma\bigl(h(Q)\bigr)^{n}\le\mu\bigl(h(Q)\bigr).
\]
由此推知
\[
\mu(\Omega)\le\mu\bigl(h(\Omega)\bigr)
\]
对每个 $\Omega\in\mathcal{O}$ 成立。事实上，设 $\Omega\in\mathcal{O}$；%
那么给定 $\varepsilon>0$，依据\trLebesgue{}\trMeasure{}的\trRegularity，%
我们可以找到含于 $\Omega$ 内的有限多个互不相交的开\trCube{} $Q_j$，%
使得 $\mu(\Omega)-\varepsilon\le\sum\mu(Q_j)$。于是由上面的\trEstimate
\begin{align*}
\mu(\Omega)-\varepsilon\ &\le\ \sum\mu\bigl(h(Q_j)\bigr)\\
&=\ \mu\Bigl(h\Bigl(\bigcup\nolimits_{j}Q_j\Bigr)\Bigr)\\
&\le\ \mu\bigl(h(\Omega)\bigr).
\end{align*}
这对每个 $\varepsilon>0$ 都成立，论断随之得证。由同一论证，%
$\mu(\Omega)\le\mu\bigl(h^{-1}(\Omega)\bigr)$，从而%
$\mu(\Omega)\le\mu\bigl(h(\Omega)\bigr)\le\mu\bigl(h^{-1}\circ h(\Omega)\bigr)
=\mu(\Omega)$，命题得证。
\end{proof}

\begin{corollary}\label{cor:homeomorphism-preserving-capacity}
若 $\R^{2n}$ 的一个\trHomeomorphism{}%
保持 $\R^{2n}$ 中\trOpenSet{}的一切\trCapacity，%
那么它也保持\trLebesgue{}\trMeasure。
\end{corollary}

迄今为止，我们的全部考察都基于一个尚未建立的\trCapacity{}之存在性。%
在下一章中，我们将借助\trHamiltonianSystem{}构造一个以动力学方式定义的%
非常特殊的\trCapacity{}函数。
